\documentclass{article}
\usepackage[T1]{fontenc}
\usepackage{lmodern}
\usepackage{graphicx}
\usepackage{amsmath,amssymb}
\usepackage[a4paper,margin=2cm]{geometry}
\usepackage[absolute,overlay]{textpos}
\setlength{\TPHorizModule}{1mm}
\setlength{\TPVertModule}{1mm}
\usepackage[indonesian]{babel}
\usepackage{hyperref}
\setlength{\emergencystretch}{2em}
% unit: Halaman 3

\begin{document}

% === Halaman 3 (python/native) ===

Sejarah DES

• Pada tahun 1972 NIST meminta standard enkripsi cipher blok.

• Pada tahun 1974 Horst Feistel di IBM mendesain cipher blok Lucifer (panjang kunci 

128 bit, Panjang blok 128 bit)

• Lucifer kemudian berkembang menjadi DES dengan beberapa masukan dari NSA 

(National Security Agency):

    - panjang kunci 56 bit, ukuran blok 64 bit

    - 16 putaran. 

• Tahun 1976 DES disetujui oleh National Bureau of Standard (NBS) setelah penilaian 

kekuatannya oleh NSA.

• Sejak itu DES diimplementasikan secara luas, baik software maupun hardware.

3

\end{document}
